\section{Interpretive Stress Tests}
\subsection{Alternative explanations are not the same as contradictions}
The observed $D=5$ is compatible with several different processes. A university can act through academic-integrity authority before a ministry publishes a cross-sector HE instrument; a ministry may communicate requirements without an accessible circular; or a binding document may be stored on an unindexed policy portal. Any of these can coexist with the locked $\stE/\stP/\stU$ evidence matrix. Distinguishing them requires data outside an official-web documentary search. Our result therefore rejects a literal claim that \emph{all four layers must be visibly explicit whenever any university rule is visible}; it does not reject coordinated government guidance, local compliance, or policy diffusion as real-world mechanisms.

Three counterfactuals clarify the observed versus non-observed margins. First, suppose a new valid, pre-cutoff HE governance control for Oman is recovered. The primary matrix and the published coding decision would need a transparent source-admission revision: Oman would change $L3:\stU\rightarrow\stE$ and $D$ would fall by one. Second, suppose an already indexed policy merely includes a new summary sentence mentioning AI but provides no control or authority. It would be $\stP$, not a valid promotion to $\stE$. Third, suppose an identified university policy explicitly cites its national strategy. That new evidence would support a specific directed-link observation $R_{1,4,j}$ but leave the existing four-layer availability vector unchanged. These experiments illustrate how the construct prevents evidence from being counted at the wrong level.

\subsection{A typed error register rather than one uncertainty score}
The sources of uncertainty are heterogeneous and cannot be collapsed into a single fuzzy membership interval. \emph{Source recovery} concerns missed qualifying instruments; \emph{source identity} concerns edition and effective date; \emph{semantic coding} concerns whether a located clause meets a control gate; \emph{language authority} concerns which version governs when bilingual texts diverge; \emph{selection-boundary error} changes panel membership; and \emph{web disappearance} affects reproducibility of a locator. These mechanisms demand different controls: bilingual repeat search, signed source snapshot, blind recoding, recorded language precedence, official eligibility ledger, and immutable archiving, respectively. The $\stP\mapsto[0,1]$ interval only represents one declared \emph{interpretation range} for already located partial evidence. It is not a sensitivity analysis over all possible missing evidence or all human judgements.

\begin{table}[t]
\centering\caption{Threat-to-evidence mapping for audit replication}\label{tab:threats}\footnotesize
\begin{tabularx}{\columnwidth}{@{}p{0.33\columnwidth}X@{}}
\toprule
Threat & Required independent control\tabularnewline\midrule
Unseen official document & Replicate Arabic/English search, register recovery\tabularnewline
Wrong edition or date & Record document version, cutoff, source hash\tabularnewline
Ambiguous control scope & Blind second coding of clause and gate\tabularnewline
Translation/authority conflict & Read controlling language; log adjudication\tabularnewline
Wrong sample boundary & Freeze regulator eligibility and founding-year list\tabularnewline
Broken/moved URL & Archive source snapshot and capture time\tabularnewline
\bottomrule
\end{tabularx}
\end{table}

\subsection{What would overturn the narrower finding?}
For the primary categorical result to be eliminated entirely, every one of the five jurisdictions currently meeting $L4=\stE,L3\ne\stE$ would have to lose its L4 explicit witness or acquire a qualifying L3 explicit HE instrument (or a combination of such changes). This is the content of $D\geq5-k$, not evidence that five independent coding errors are improbable. If only one decisive record is disputed, the matrix difference is at most one jurisdictional indicator in this specific primary result. This claim remains conditional on the six official-source profiles and does not prevent a later, better documented corpus from substantially revising the map.

The broader UAE-OBEF interpretation is informative because it deliberately expands what may count as sector translation but still leaves three L3 unlocated jurisdictions. It is \emph{not} a test of implementation, and the fact that the partial-state interval upper bound stays fixed does not show that the policy content is invariant. The two kinds of sensitivity answer separate questions: alternative L3 gate definitions change a nominal judgement; the interval envelope allows any membership within partial cases given the observed judgement. Reporting both and calling neither a sampling confidence interval makes the findings more reusable across new evidence windows.
